% from ega4-errata-addenda-list3-fr.tex
\oldpage[IV]{348}
\paragraph{(Err\textsubscript{IV}, 3).}
Dans (III, 4.4.9), ligne 15 du bas, p.~137, remplacer : finies par :
finies en tant qu'ensembles.

\paragraph{(Err\textsubscript{IV}, 4).}
Dans (III, 7.8.10), ligne 9 de la p.~76, remplacer simple par lisse.

\paragraph{(Err\textsubscript{IV}, 5).}
Dans ($0_{\mathrm{IV}}$, 16.1.4), ligne 8 du bas de la p.~23, remplacer :
Dire que $A$ est caténaire, par : Dire qu'un anneau intègre $A$ est
caténaire.

\paragraph{(Err\textsubscript{IV}, 6).}
Dans ($0_{\mathrm{IV}}$, 16.1.9), on peut supprimer l'hypothèse que
$A$ est noethérien. En effet, pour tout élément $u\in B_M$, $A[u]$
est une sous-$A$-algèbre de $C$, donc $u$ est entier sur $A$ (Bourbaki,
\emph{Alg. comm.}, chap.~V, \S1, n\textsuperscript{o}~1, cor. de la
prop.~1), et par suite aussi sur $A_M$; on conclut encore par
{\hyperlink{ega4.c04.prop.16-1-5}{(16.1.5)}}.

\paragraph{(Err\textsubscript{IV}, 7).}
Dans ($0_{\mathrm{IV}}$, 16.5.12), on notera que la relation
(16.5.11.1) équivaut à la même relation où l'on remplace $A$ par les
anneaux quotients $A/\mathfrak r$, $\mathfrak r$ parcourant l'ensemble
des idéaux premiers minimaux de $A$. On peut donc dans la preuve de
{\hyperlink{ega4.c04.statement.16-5-12}{(16.5.12)}} se borner au cas
où $A$ est intègre, et appliquer (16.1.4.2).

\paragraph{(Err\textsubscript{IV}, 8).}
Supprimer ($0_{\mathrm{IV}}$, 19.3.12) dont la démonstration est erronée
(les ensembles $W_\lambda$, $U_\lambda$ et $V_\lambda$ ne sont pas des
variétés linéaires).

\paragraph{(Err\textsubscript{IV}, 9).}
Dans ($0_{\mathrm{IV}}$, 19.8.3), lignes 6 et 8 du bas de la p.~110,
remplacer homomorphisme par homomorphisme local; de même dans
($0_{\mathrm{IV}}$, 19.8.4), ligne 1 de la p.~111.

\paragraph{(Err\textsubscript{IV}, 10).}
Dans ($0_{\mathrm{IV}}$, 20.1.2), ligne 2 du bas de la p.~117,
remplacer $a\in A$ par $b\in B$ et $aD$ par $bD$; ligne 1 du bas de la
p.~117, remplacer $A$-module par $B$-module.

\paragraph{(Err\textsubscript{IV}, 11).}
Dans ($0_{\mathrm{IV}}$, 20.2.1), ligne 8 de la p.~119, remplacer
$A$-modules par $B$-modules. Lignes 21 et 22 de la p.~119, remplacer
« di-homomorphismes de modules (relatif à $u$) » par : homomorphisme
de $B$-modules.

\paragraph{(Err\textsubscript{IV}, 12).}
Dans ($0_{\mathrm{IV}}$, 20.4.6.1), remplacer $p_2-p_1$ par $p_1-p_2$;
ligne 10 du bas de la p.~125, remplacer $x\otimes1-1\otimes x$ par
$1\otimes x-x\otimes1$.

\paragraph{(Err\textsubscript{IV}, 13).}
Dans ($0_{\mathrm{IV}}$, 20.4.8), ligne 14 de la p.~126, remplacer
$p_2=p_1-d_{B/A}$ par $p_1=p_2-d_{B/A}$.

\paragraph{(Err\textsubscript{IV}, 14).}
Dans ($0_{\mathrm{IV}}$, 20.4.13), remplacer les lignes 6 à 11 du bas
de la p.~127 par le texte suivant :

\medskip
\noindent
(iv) Soient $A$, $B$ deux anneaux intègres tels que $A\subset B$, $B$
étant entier sur $A$, et que le corps des fractions de $B$ soit une
extension séparable de celui de $A$. Alors le $B$-module
$\Omega^1_{B/A}$ est un module de torsion. En effet, pour tout $x\in B$,
soit $f$ le polynôme minimal de $x$ par rapport au corps des fractions
$K$ de $A$; il existe $a\in A$ non nul tel que $af[T]\in A[T]$; comme
$x$ est séparable sur $K$, on a $f'(x)\neq0$, et d'autre part de la
relation $af(x)=0$ on déduit $af'(x)d_{B/A}x=0$, d'où notre assertion
en vertu de {\hyperlink{ega4.c04.cor.20-4-7}{(20.4.7)}}.

\paragraph{(Err\textsubscript{IV}, 15).}
Dans ($0_{\mathrm{IV}}$, 22.2.7), ligne 5 de la p.~189, au lieu de
« formellement lisses », il faut « formellement lisses pour les
topologies préadiques ».

\paragraph{(Err\textsubscript{IV}, 16).}
Dans le Sommaire du chap.~IV, ligne 4 du bas de la p.~222, ajouter :
et anneaux strictement locaux. Lignes 1 à 3 du bas, remplacer par :
\begin{quote}
\S19. Immersions régulières et platitude normale.\\
\S20. Fonctions méromorphes et pseudo-morphismes.\\
\S21. Diviseurs.
\end{quote}
Ligne 16 de la p.~223, remplacer « de présentation finie » par
« localement de présentation finie »; ligne 19 de la p.~223, remplacer
« comme » par « comme étant ».

\paragraph{(Err\textsubscript{IV}, 17).}
Dans (IV, 3.1.10), ajouter après la ligne 16 de la p.~39 :

\begin{remark}[3.1.10.2]
\label{IV.3.1.10.2}
Soit $f:X\to Y$ un morphisme quasi-compact et quasi-séparé, et soit
$\sh{F}$ un $\sh{O}_X$-Module quasi-cohérent, de sorte que
$f_*(\sh{F})$ est un $\sh{O}_Y$-Module quasi-cohérent
{\hyperref[canon-index.ega4.c04.statement.1-7-4]{(1.7.4)}}. Alors on a
\begin{align}
 \operatorname{Ass}(f_*(\sh{F}))&\subset f(X),
 \label{IV.3.1.10.3}\tag{3.1.10.3}\\
 \operatorname{Ass}(f_*(\sh{F}))&\subset f(\operatorname{Ass}(\sh{F}))
 \quad\text{si $X$ est localement noethérien}.
 \label{IV.3.1.10.4}\tag{3.1.10.4}
\end{align}

Soit en effet $y\in\operatorname{Ass}(f_*(\sh{F}))$; en vertu de
{\hyperref[canon-index.ega4.c04.lem.2-3-1]{(2.3.1)}} et du fait que le morphisme
$\operatorname{Spec}(\sh{O}_{Y,y})\to Y$ est plat,
$(f_*(\sh{F}))_y$ est égal à $(f'_*(\sh{F}'))_y$, où
\[
 f':X\times_Y\operatorname{Spec}(\sh{O}_{Y,y})
 \longrightarrow\operatorname{Spec}(\sh{O}_{Y,y})
\]
est déduit de $f$ par changement de base et
$\sh{F}'=\sh{F}\otimes_{\sh{O}_Y}\sh{O}_{Y,y}$. On est donc ramené à
prouver (3.1.10.3) et (3.1.10.4) en y remplaçant $Y$ par
$\operatorname{Spec}(\sh{O}_{Y,y})$, $f$ par $f'$ et $\sh{F}$ par
$\sh{F}'$, compte tenu de (I, 3.6.5) et de la définition de
$\operatorname{Ass}(\sh{F}')$.

Par définition, il y a un élément
$t_y\in(f_*(\sh{F}))_y$ non nul et dont $\mathfrak m_y$ est
l'annulateur, et comme $(f_*(\sh{F}))_y$ s'identifie à
$\Gamma(Y,f_*(\sh{F}))$ puisque $Y$ est local, $t_y$ s'identifie aussi
à une section $s\in\Gamma(X,\sh{F})$. L'Idéal de $\sh{O}_X$ annulateur
de $s$ contient donc $\mathfrak m_y\sh{O}_X$, et il résulte de
($0_{\mathrm I}$, 1.7.4) que le support de $s$ est contenu dans
$f^{-1}(y)$; si ce dernier était vide, on aurait $s=0$, ce qui est
absurde. Lorsque de plus $X$ est localement noethérien, alors le
sous-$\sh{O}_X$-Module cohérent $\sh{O}_Xs$ de $\sh{F}$ n'est pas nul,
donc $\operatorname{Ass}(\sh{O}_Xs)\neq\varnothing$
{\hyperref[canon-index.ega4.c04.cor.3-1-5]{(3.1.5)}}; mais
\[
 \operatorname{Ass}(\sh{O}_Xs)\subset
 \operatorname{Supp}(\sh{O}_Xs)\subset f^{-1}(y),\qquad
 \operatorname{Ass}(\sh{O}_Xs)\subset\operatorname{Ass}(\sh{F})
\]
{\hyperref[canon-index.ega4.c04.prop.3-1-7]{(3.1.7)}}, donc
$y\in f(\operatorname{Ass}(\sh{F}))$.
\end{remark}

